
    %%%%%%%%%%%%%%%%%%%%%%%%%%%%%%%%%%%%%%%%%%%%%%%%%%%%%%%%%%%%%%%%%%%%%%%%%%%%%%%%
    %%% Classe do documento
    \documentclass[12pt, addpoints]{exam}
    \UseRawInputEncoding

    %%%%%%%%%%%%%%%%%%%%%%%%%%%%%%%%%%%%%%%%%%%%%%%%%%%%%%%%%%%%%%%%%%%%%%%%%%%%%%%%
    % preambulo.tex
    %
    % Arquivo de configuração de packages para uso o LaTeX, conforme minhas
    % preferências e modelos pessoais.
    %
    % Para maiores informações, visite:
    %    https://github.com/abrantesasf/latex
    %
    % NÃO ALTERE SE NÃO SOUBER O QUE ESTÁ FAZENDO!
    %%%%%%%%%%%%%%%%%%%%%%%%%%%%%%%%%%%%%%%%%%%%%%%%%%%%%%%%%%%%%%%%%%%%%%%%%%%%%%%%


    %%%%%%%%%%%%%%%%%%%%%%%%%%%%%%%%%%%%%%%%%%%%%%%%%%%%%%%%%%%%%%%%%%%%%%%%%%%%%%%%
    %%% Estruturas de controle:
    \input{static/utils/controles}

    %%%%%%%%%%%%%%%%%%%%%%%%%%%%%%%%%%%%%%%%%%%%%%%%%%%%%%%%%%%%%%%%%%%%%%%%%%%%%%%%
    %%% Compilação condicional em PDF:
    \input{static/utils/pdf}

    %%%%%%%%%%%%%%%%%%%%%%%%%%%%%%%%%%%%%%%%%%%%%%%%%%%%%%%%%%%%%%%%%%%%%%%%%%%%%%%%
    %%% Configurações de layout da página:
    \input{static/utils/layout}

    %%%%%%%%%%%%%%%%%%%%%%%%%%%%%%%%%%%%%%%%%%%%%%%%%%%%%%%%%%%%%%%%%%%%%%%%%%%%%%%%
    %%% Configurações para autores:
    \input{static/utils/autores}

    %%%%%%%%%%%%%%%%%%%%%%%%%%%%%%%%%%%%%%%%%%%%%%%%%%%%%%%%%%%%%%%%%%%%%%%%%%%%%%%%
    %%% Cabeçalho e rodapé:
    %%% Atenção! É MUITO DIFÍCIL ajustar apropriadamente os running headers
    %%% e running footers da primeira vez. Você pode confiar no LaTeX e deixar que
    %%% ele faça o ajuste automaticamente ou pode quebrar a cabeça e
    %%% tentar melhorar algo que o LaTeX já faz muito bem, mesmo que não fique
    %%% exatamente do jeito que você quer. O que você prefere? Se quiser ajustar
    %%% manualmente, descomente a linha abaixo e faça as configurações no arquivo
    %%% "utils/headings.tex".
    %\input{static/utils/headings}

    %%%%%%%%%%%%%%%%%%%%%%%%%%%%%%%%%%%%%%%%%%%%%%%%%%%%%%%%%%%%%%%%%%%%%%%%%%%%%%%%
    %%% Configuração para as fontes do documento:
    %%% Se você quiser usar fontes próprias ou do sistema, precisará ajustar as
    %%% configurações no arquivo "utils/fontes.tex".
    %%% ATENÇÃO: por padrão eu utilizo XeLaTeX com fontes proprietárias projetadas
    %%% por Matthew Butterick, adquiridas em https://mbtype.com, que não podem ser
    %%% distribuídas devido à licença de utilização. Se você usar XeLaTeX, você
    %%% DEVERÁ OBRIGATORIAMENTE ajustar as fontes no arquivo "utils/fontes.tex".
    \input{static/utils/fontes}

    %%%%%%%%%%%%%%%%%%%%%%%%%%%%%%%%%%%%%%%%%%%%%%%%%%%%%%%%%%%%%%%%%%%%%%%%%%%%%%%%
    %%% Sumário:
    \input{static/utils/toc}

    %%%%%%%%%%%%%%%%%%%%%%%%%%%%%%%%%%%%%%%%%%%%%%%%%%%%%%%%%%%%%%%%%%%%%%%%%%%%%%%%
    %%% Matemática:
    \input{static/utils/matematica}

    %%%%%%%%%%%%%%%%%%%%%%%%%%%%%%%%%%%%%%%%%%%%%%%%%%%%%%%%%%%%%%%%%%%%%%%%%%%%%%%%
    %%% Notas de rodapé:
    \input{static/utils/footnotes}

    %%%%%%%%%%%%%%%%%%%%%%%%%%%%%%%%%%%%%%%%%%%%%%%%%%%%%%%%%%%%%%%%%%%%%%%%%%%%%%%%
    %%% Referências cruzadas, links, citações:
    \input{static/utils/crossrefs}

    %%%%%%%%%%%%%%%%%%%%%%%%%%%%%%%%%%%%%%%%%%%%%%%%%%%%%%%%%%%%%%%%%%%%%%%%%%%%%%%%
    %%% Configurações para os hiperlinks em PDF:
    \input{static/utils/pdfconfig}

    %%%%%%%%%%%%%%%%%%%%%%%%%%%%%%%%%%%%%%%%%%%%%%%%%%%%%%%%%%%%%%%%%%%%%%%%%%%%%%%%
    %%% Glossário e índice remissivo:
    \input{static/utils/glossindex}

    %%%%%%%%%%%%%%%%%%%%%%%%%%%%%%%%%%%%%%%%%%%%%%%%%%%%%%%%%%%%%%%%%%%%%%%%%%%%%%%%
    %%% Referências bibliográficas:
    \input{static/utils/bibliografia}

    %%%%%%%%%%%%%%%%%%%%%%%%%%%%%%%%%%%%%%%%%%%%%%%%%%%%%%%%%%%%%%%%%%%%%%%%%%%%%%%%
    %%% Cores:
    \input{static/utils/cores}

    %%%%%%%%%%%%%%%%%%%%%%%%%%%%%%%%%%%%%%%%%%%%%%%%%%%%%%%%%%%%%%%%%%%%%%%%%%%%%%%%
    %%% Computação:
    \input{static/utils/computacao}

    %%%%%%%%%%%%%%%%%%%%%%%%%%%%%%%%%%%%%%%%%%%%%%%%%%%%%%%%%%%%%%%%%%%%%%%%%%%%%%%%
    %%% Gráficos:
    \input{static/utils/graficos}

    %%%%%%%%%%%%%%%%%%%%%%%%%%%%%%%%%%%%%%%%%%%%%%%%%%%%%%%%%%%%%%%%%%%%%%%%%%%%%%%%
    %%% Ícones:
    \input{static/utils/icones}

    %%%%%%%%%%%%%%%%%%%%%%%%%%%%%%%%%%%%%%%%%%%%%%%%%%%%%%%%%%%%%%%%%%%%%%%%%%%%%%%%
    %%% Blocos:
    \input{static/utils/blocos}

    %%%%%%%%%%%%%%%%%%%%%%%%%%%%%%%%%%%%%%%%%%%%%%%%%%%%%%%%%%%%%%%%%%%%%%%%%%%%%%%%
    %%% Boxes coloridos:
    \input{static/utils/colorbox}

    %%%%%%%%%%%%%%%%%%%%%%%%%%%%%%%%%%%%%%%%%%%%%%%%%%%%%%%%%%%%%%%%%%%%%%%%%%%%%%%%
    %%% Tabelas:
    \input{static/utils/tabelas}

    %%%%%%%%%%%%%%%%%%%%%%%%%%%%%%%%%%%%%%%%%%%%%%%%%%%%%%%%%%%%%%%%%%%%%%%%%%%%%%%%
    %%% Ambientes de listas:
    \input{static/utils/listas}

    %%%%%%%%%%%%%%%%%%%%%%%%%%%%%%%%%%%%%%%%%%%%%%%%%%%%%%%%%%%%%%%%%%%%%%%%%%%%%%%%
    %%% Ambientes floats e similares:
    \input{static/utils/floats}

    %%%%%%%%%%%%%%%%%%%%%%%%%%%%%%%%%%%%%%%%%%%%%%%%%%%%%%%%%%%%%%%%%%%%%%%%%%%%%%%%
    %%% Ajustes para a classe EXAM:
    \input{static/utils/exam}

    %%%%%%%%%%%%%%%%%%%%%%%%%%%%%%%%%%%%%%%%%%%%%%%%%%%%%%%%%%%%%%%%%%%%%%%%%%%%%%%%
    %%% Ajustes para textos:
    \input{static/utils/texto}

    %%%%%%%%%%%%%%%%%%%%%%%%%%%%%%%%%%%%%%%%%%%%%%%%%%%%%%%%%%%%%%%%%%%%%%%%%%%%%%%%
    %%% Ambientes, aliases, comandos e customizações em geral para serem
    %%% utilizadas nos documentos. Defina aqui qualquer customização específica
    %%% para seus documentos!
    \input{static/utils/customizacoes}

    %%%%%%%%%%%%%%%%%%%%%%%%%%%%%%%%%%%%%%%%%%%%%%%%%%%%%%%%%%%%%%%%%%%%%%%%%%%%%%%%
    %%% Ajuste do layout, espaçamento de linhas e etc.:
    \geometry{a4paper, portrait, left=2.0cm, right=2.0cm, top=2.0cm, bottom=2.0cm}
    \singlespacing

    %%%%%%%%%%%%%%%%%%%%%%%%%%%%%%%%%%%%%%%%%%%%%%%%%%%%%%%%%%%%%%%%%%%%%%%%%%%%%%%%
    %%% Configurações de cabeçalho e rodapé (não use fancyhdr pois ocorre conflito):
    \pagestyle{headandfoot}
    \runningheadrule
    \firstpageheader{}{}{}
    \runningheader{Sem disciplina}{}{Avaliação}
    \firstpagefooter{}{}{Página \thepage\ de \numpages}
    \runningfooter{}{}{Página \thepage\ de \numpages}
    \runningfootrule

    %%%%%%%%%%%%%%%%%%%%%%%%%%%%%%%%%%%%%%%%%%%%%%%%%%%%%%%%%%%%%%%%%%%%%%%%%%%%%%%%
    %%% Configurações para as propriedades do PDF:
    \hypersetup{
    hidelinks,           % Comente para web, descomente para impressão
    colorlinks=true,      % True para web, False para impressão
    pdftitle=50: Avaliação,
    pdfauthor=Matheus Endlich Silveira,
    pdfsubject=Sem disciplina,
    pdfkeywords=['Sem tag'],
    pdfinfo={
        CreationDate={}, % Ex.: D:AAAAMMDDHH24MISS
        ModDate={}       % Ex.: D:AAAAMMDDHH24MISS
    }
    }
    \input{static/utils/pdfconfig}

    %%%%%%%%%%%%%%%%%%%%%%%%%%%%%%%%%%%%%%%%%%%%%%%%%%%%%%%%%%%%%%%%%%%%%%%%%%%%%%%%
    %%% Começa o documento
    \begin{document}

    %%%%%%%%%%%%%%%%%%%%%%%%%%%%%%%%%%%%%%%%%%%%%%%%%%%%%%%%%%%%%%%%%%%%%%%%%%%%%%%%
    %%% Folha de instruções padronizadas da UVV para a prova
    \pdfbookmark[1]{Instruções}{instrucoes}
    \begin{figure}[H]
        \begin{center}
            \includegraphics[scale=0.3]{{static/imgs/logo-uvv.png}}
        \end{center}	
    \end{figure}

    \vspace{-1cm}

    \begin{table}[H]
        \centering
        \begin{tabular}{|llll|}
            \hline
            \multicolumn{2}{|l|}{\textbf{Disciplina:} Sem disciplina} & 
            \multicolumn{1}{l|}{\multirow{3}{*}{\textbf{\makecell{Nota:\\ \,\\ \,}}}} & 
            \multirow{3}{*}{\textbf{\makecell{Coordenador:\\ \,\\ \,}}} \\ 
            \cline{1-2}
            \multicolumn{2}{|l|}{\textbf{Professor:} Matheus Endlich Silveira \hspace{4.72cm} } & 
            \multicolumn{1}{l|}{} & \\ 
            \cline{1-2}
            \multicolumn{2}{|l|}{\textbf{Aluno:}} & 
            \multicolumn{1}{l|}{} & \\ 
            \hline
            \multicolumn{1}{|l|}{\textbf{Turma:} \hspace{6.3cm} } & 
            \multicolumn{1}{l|}{\textbf{Semestre:}} & 
            \multicolumn{2}{l|}{\textbf{Valor:} 10 pontos} \\ 
            \hline
            \multicolumn{1}{|l|}{\textbf{Data:}} & 
            \multicolumn{3}{l|}{\textbf{Avaliação:}} \\ 
            \hline
        \end{tabular}
    \end{table}

    \begin{center}
        \fbox{\setlength\fboxsep{0.3cm} \fbox{\parbox{15.4cm}{
            \begin{center}
                \textbf{INSTRUÇÕES PARA A REALIZAÇÃO DESTA AVALIAÇÃO:}
            \end{center}
            \begin{itemize}[leftmargin=*]
                \item Esta é a primeira avaliação da disciplina Sem disciplina, 
                    e engloba o conteúdo referente Sem tag.
                \item Esta avaliação contém 20 questões objetivas, \textbf{todas obrigatórias}.
                \item \textbf{Leia atentamente cada questão!} As questões objetivas terão uma,
                    e apenas uma única, resposta a ser marcada.
                \item \textbf{Desligue o celular} antes de começar. A avaliação é
                    \textbf{individual} e \textbf{sem consulta}.
                \item As questões podem ser respondidas, na prova, com lápis ou caneta. Ao final
                    da prova \textbf{{VOCÊ DEVERÁ PREENCHER O GABARITO COM CANETA
                    DE COR PRETA}} \textbf{\underline{OU AZUL}} para que seja feita a correção
                    automatizada através da leitura do padrão de respostas marcado no
                    gabarito.
                \item \textbf{CUIDADO ao preencher o gabarito} pois eles são \textbf{INDIVIDUAIS
                    e NOMEADOS} (cada estudante tem um gabarito próprio específico, com um
                    código de identificação único). \textbf{Questões rasuradas são anuladas}
                    pelo software de leitura do gabarito.
                \item Ao final da prova, \textbf{DEVOLVA A PROVA E O GABARITO} para o professor.
                \item Siga todas as normas de \textbf{Integridade Acadêmica} da disciplina.
                    Alunos que forem flagrados com qualquer espécie de ``cola'' ou trocando
                    informações com outros alunos terão suas avaliações recolhidas, as notas
                    zeradas, e a situação será encaminhada para a coordenação para a aplicação
                    das penalidades previstas pela universidade.
                \item Com 90 minutos de avaliação você terá tempo de até 4,min por questão,
                    em média.
                \item Boa avaliação!
            \end{itemize}
            \vspace{2cm}
        }}}
    \end{center}
    \newpage

    %%%%%%%%%%%%%%%%%%%%%%%%%%%%%%%%%%%%%%%%%%%%%%%%%%%%%%%%%%%%%%%%%%%%%%%%%%%%%%%%
    %%% Inicia o bloco de questões:
    \begin{questions}

    \newpage
    \question

Quais das seguintes afirmações são verdadeiras? As Métricas de software servem para:



\begin{itemize}

    \item[I)] indicar a qualidade do produto e avaliar a produtividade.

    \item[II)] auxiliar na melhoria do processo.

    \item[III)] formar uma base para as estimativas e justificar a aquisição de ferramentas.

    \item[IV)] determinar se a utilização de um método traz benefícios ou não.

\end{itemize}
\begin{choices}
\choice Apenas as alternativas II e III.
\choice Todas as alternativas.
\choice Nenhuma delas.
\choice Apenas as alternativas I, IV.
\choice Apenas as alternativas I, II e IV.
\end{choices}

\question

O usuário que não é da área de tecnologia, costuma ser de alta gerência, que é

especialista no negócio e tem a responsabilidade central pelas políticas de

dados da empresa é o:
\begin{choices}
\choice Coordenador de Tecnologia da Informação (ITS: IT Supervisor)
\choice Administrador do Banco de Dados (DBA: Database Administrator)
\choice Administrador de Dados (DA: Data Administrator)
\choice Encarregado da Proteção de Dados (DPO: Data Protection Officer)
\choice Gerente de Tecnologia da Informação (ITM: IT Manager)
\end{choices}

\question

Seu chefe que fazer uma campanha de venda para produtos que custam entre 20 e 30

reais, mas está preocupado porque ouviu de um dos vendedores que diversos

produtos nessa faixa de preço estão com o estoque zerado. Para verificar a

situação seu chefe pediu para que você prepare um relatório que mostre os

produtos que custam entre 20 e 30 reais e que estão com o estoque zerado. Você

preparou o seguinte relatório para seu chefe:

\begin{tcolorbox}

 \begin{verbatim}1  SELECT l.nome AS loja

2       , p.nome AS produto

3       , p.preco_unitario AS preco

4       , e.quantidade AS qtd

5  FROM produtos p

6  INNER JOIN estoques e ON (e.produto_id = p.produto_id)

7  INNER JOIN lojas l ON (e.loja_id = l.loja_id)

8  WHERE (p.preco_unitario >= 20 OR p.preco_unitario <= 30)

9    AND e.quantidade = 0

10 ORDER BY loja, produto;

 \end{verbatim}

\end{tcolorbox}

O seu relatório:
\begin{choices}
\choice Está errado pois a tabela que deveria estar na cláusula FROM é a tabela ``lojas', não a tabela ``produtos'.
\choice Está errado pois vai trazer também produtos fora da faixa de preço que o seu chefe quer.
\choice Nenhuma das alternativas anteriores.
\choice Está correto e trará exatamente, os produtos com o estoque zerado e na faixa de preço requisitada pelo seu chefe. Além disso ordena o resultado pelo nome da loja e pelo nome do produto.
\choice Está errado pois há um erro de sintaxe nas linhas 6 e 7.
\end{choices}

\question

Supondo a Relação $\texttt{PROJ (PNO, Orçam)}$, com chave primária $\texttt{PNO}$, a Relação $\texttt{EMP (ENO, ENome, Cargo)}$ com chave primária $\texttt{ENO}$, e a Relação $\texttt{DSG (ENO, PNO, Dur, Resp)}$, com chave primária \{\texttt{ENO, PNO}\}, chave estrangeira $\texttt{PNO}$ em relação a $\texttt{PROJ}$ e chave estrangeira $\texttt{ENO}$ em relação a $\texttt{EMP}$. Qual das expressões da álgebra relacional abaixo $\textbf{NÃO}$ corresponde à seguinte consulta SQL:



\begin{verbatim}

SELECT ENome

FROM EMP, PROJ, DSG

WHERE EMP.ENO = DSG.ENO

      AND PROJ.PNO = DSG.PNO

      AND Dur > 36

\end{verbatim}
\begin{choices}
\choice \( \pi_{ENome} (\sigma_{Dur > 36} ((\pi_{PNO} (\text{PROJ})) \Join_{PNO} (\text{EMP} \Join_{ENO} (\pi_{Dur} (\text{DSG})))) \)
\choice \( \pi_{ENome} (\text{PROJ} \Join_{PNO} (\text{EMP} \Join_{ENO} \sigma_{Dur > 36} (\pi_{Dur} (\text{DSG})))) \)
\choice \( \pi_{ENome} (\sigma_{Dur > 36} ((\pi_{ENome, ENO} (\text{EMP})) \Join_{ENO} (\text{PROJ} \Join_{PNO} (\text{DSG})))) \)
\choice \( \pi_{ENome} (\text{PROJ} \Join_{PNO} (\text{EMP} \Join_{ENO} \sigma_{Dur > 36} (\text{DSG}))) \)
\choice \( \pi_{ENome} (\text{PROJ} \Join_{PNO} (\sigma_{Dur > 36} (\text{EMP} \Join_{ENO} (\text{DSG})))) \)
\end{choices}

\question

Numa prova de múltipla escolha com 10 questões e 4 alternativas, qual a chance (probabilidade) de um aluno apenas "chutando as respostas" conseguir "gabaritar" a prova (acertar todas as questões).


\begin{choices}
\choice \( \frac{1}{10^{8}} \)
\choice \( \frac{1}{2^{20}} \)
\choice \( \frac{1}{10^{4}} \)
\choice \( \frac{1}{4^{20}} \)
\choice \( \frac{1}{4^{15}} \)
\end{choices}

\question

Em uma lista circular duplamente encadeada com \( n \) elementos, o espaço ocupado apenas pelos apontadores é (assuma que um apontador ocupa \( p \) bytes):
\begin{choices}
\choice \( (np)^2 \)
\choice \( np \)
\choice \( 6np \)
\choice \( 4np \)
\choice \( 2np \)
\end{choices}

\question

Quais dos algoritmos de ordenação abaixo possuem tempo no pior caso e tempo médio de execução proporcional a \( O(n \log n) \).
\begin{choices}
\choice Quicksort e merge sort
\choice Bubble sort e Quick sort
\choice Merge sort e bubble sort
\choice Heap sort e selection sort
\choice Merge sort e heap sort
\end{choices}

\question

Em relação aos metadados de um banco de dados, assinale a afirmação correta:
\begin{choices}
\choice São atualizados pelo SGBD, sempre que ocorre alguma operação de inserção, atualização ou remoção de dados.
\choice Armazenam informações como os dados de cadastro dos usuários, os valores de compra em uma ordem de compra, as turmas e disciplinas que os professores dão aula, e dados de recursos humanos.
\choice Como são sempre atualizados de forma automática pelo sistema de gerenciamento de bancos de dados, estão sempre prontos para nos dar uma ``foto' da estrutura do banco de dados em um momento no tempo.
\choice São armazenados no catálogo (ou dicionário de dados) do sistema de gerenciamento de bancos de dados, que fica nas mesmas tabelas de dados dos usuários, permitindo assim que os metadados estejam sempre atualizados.
\choice Armazenam informações sobre os dados dos usuários, ou seja, são dados de rotina dos usuários, aqueles dados que os usuários precisam para trabalhar no seu dia a dia.
\end{choices}

\question

\textbf Supondo a Relação \texttt{PROJ (PNO, Nome, Orçam)}, com chave primária \texttt{PNO} e a Relação \texttt{DSG (ENO, PNO, Dur, Resp)}, com chave primária \{\texttt{ENO, PNO}\} e chave estrangeira \texttt{PNO} em relação a \texttt{PROJ}, a asserção abaixo \textbf{NÃO} expressa:



\begin{equation}\forall g \in \texttt{DSG}, \exists j \in \texttt{PROJ} : g.\texttt{PNO} = j.\texttt{PNO}\end{equation}
\begin{choices}
\choice Uma restrição a ser verificada na atualização de tuplas em $\texttt{DSG}$.
\choice Uma restrição a ser verificada na inserção de tuplas em $\texttt{DSG}$.
\choice Uma restrição de integridade de chave primária em $\texttt{PROJ}$.
\choice Uma restrição de integridade de chave estrangeira em $\texttt{DSG}$.
\choice Uma restrição que define um estado consistente do banco de dados.
\end{choices}

\question

Uma característica de uma arquitetura RISC é:
\begin{choices}
\choice Possui um pequeno conjunto de instruções simples
\choice A arquitetura é constituída de milhares de processadores
\choice Uma arquitetura de alto risco pois o mercado de hardware evolui muito rapidamente
\choice Possui um grande conjunto de instruções complexas
\choice O processador é formado por válvulas e transistores
\end{choices}

\question

A sentença descritiva resultante de um processo de mensuração, correspondendo a

um fato bruto observado/conhecido que pode ser registrado, que não foi

processado para revelar seu significado/importância, que tem um certo

significado implícito e que deve estar corretamente formatado para o

armazenamento, processamento e apresentação é:
\begin{choices}
\choice Informação
\choice Processamento
\choice Dado
\choice Conhecimento
\choice Sabedoria
\end{choices}

\question

Se \( \lim_{n \to \infty} \left( \frac{1}{n^2} + \frac{2}{n^2} + \cdots + \frac{n-1}{n^2} \right) = L \) então


\begin{choices}
\choice \( L = 0 \)
\choice \( L = 2 \)
\choice \( L = \infty \)
\choice \( L = 1 \)
\choice \( L = 1/2 \)
\end{choices}

\question

Uma árvore binária é declarada em C como



\begin{verbatim}

typedef struct no *apontador;

struct no {

    int valor;

    apontador esq, dir;

};

\end{verbatim}



onde \texttt{esq} e \texttt{dir} representam ligações para os filhos esquerdo e direito de um nó da árvore, respectivamente. Qual das seguintes alternativas é uma implementação correta da operação que inverte as posições dos filhos esquerdo e direito de um nó \texttt{p} da árvore, onde \texttt{t} é um apontador auxiliar.


\begin{choices}
\choice \texttt{t = p;} \\

          \texttt{p->esq = p->dir;} \\

          \texttt{p->dir = p->esq;}
\choice \texttt{p->esq = p->dir;} \\

          \texttt{t = p->esq;} \\

          \texttt{p->dir = t;}
\choice \texttt{p->dir = t;} \\

          \texttt{p->esq = p->dir;} \\

          \texttt{p->dir = t;}
\choice \texttt{t = p->dir;} \\

          \texttt{p->dir = p->esq;} \\

          \texttt{p->esq = t;}
\choice \texttt{t = p->dir;} \\

          \texttt{p->esq = p->dir;} \\

          \texttt{p->dir = t;}
\end{choices}

\question

Qual das seguintes afirmações sobre crescimento assintótico de funções não é verdadeira:
\begin{choices}
\choice Se \( f(n) = O(g(n)) \) então \( g(n) = O(f(n)) \)
\choice \( 2n^2 + 3n + 1 = O(n^2) \)
\choice \( 2^{n+1} = O(2^n) \)
\choice Se \( f(n) = O(g(n)) \) e \( g(n) = O(h(n)) \) então \( f(n) = O(h(n)) \)
\choice \( \log n^2 = O(\log n) \)
\end{choices}

\question

Quando Edgar Codd cricou o modelo relacional de dados,

também criou algumas regras para que a estrutura das tabelas fosse adequada às

operações de inserção, atualização, remoção e busca de dados. Ao avaliar um

banco de dados que armazena dados dos eleitores, dos partidos e dos candidatos

de uma eleição, você encontrou a situação ilustrada na tabela abaixo:

\begin{figure}[H]

  \centering

  \includegraphics[width=0.5\linewidth]{static/imgs/defaul.png}

\end{figure}

Em relação às boas normas de projeto do modelo relacional, assinale a

alternativa que corretamente identifica um problema no projeto desta tabela:
\begin{choices}
\choice A tabela contém um atributo multivalorado (o correto seria que cada tabela não tivesse atributos multivalorados).
\choice Apesar do telefone ser um atributo multivalorado, do jeito que a tabela está criada só é possível armazenar dois números de telefones, criando uma restrição ao armazenamento de telefones de contado de eleitores que tenham mais de dois telefones.
\choice Todas as afirmativas anteriores.
\choice A tabela está misturando dados de três entidades diferentes, eleitores, partidos e candidatos (o correto seria que cada tabela armazenasse dados de um único assunto ou entidade).
\choice A chave primária para a tabela não está definida e, somente olhando as colunas, não é possível dizer se é a coluna 'numero' ou a coluna 'titulo'.
\end{choices}

\question

Considere o seguinte trecho de programa:\\

1. \( i := 1; \)\\

2. while \( i \leq n \) do\\

   begin\\

3. \( sum := sum + a[i]; \)\\

4. \( i := i + 1; \)\\

   end;\\



Considere que:\\

- \( I \) representa a inicialização da variável \( i := 1 \) na linha 1;\\

- \( T \) representa o teste da linha 2;\\

- \( A \) representa os comandos da linha 3;\\

- \( P \) representa o incremento na linha 4.\\



Qual é a expressão regular que representa todas as sequências de passos possíveis de serem executados por este trecho de programa?
\begin{choices}
\choice \( I(T(AP)^*)^* \)
\choice \( IT^+A^*P^* \)
\choice \( I(TAP)^+ \)
\choice \( I(TAP)^* \)
\choice \( I(T(AP)^*)^+ \)
\end{choices}

\question

Considere as seguintes afirmações sobre resolução de problemas em IA.

\begin{enumerate}[label=\Roman*.]

    \item A* é um conhecido algoritmo de busca heurística.

    \item O Minimax é um dos principais algoritmos para jogos de dois jogadores, como o xadrez.

    \item Busca em espaço de estados é uma das formas de resolução de problemas em IA.

\end{enumerate}

São corretas:
\begin{choices}
\choice Apenas II e III
\choice I, II e III
\choice Apenas I e II
\choice Apenas III
\choice Apenas I e III
\end{choices}

\question

Qual das seguintes opções \textbf{não é} uma vantagem do uso de um banco de

dados sobre o sistema de arquivos convencional?


\begin{choices}
\choice Estruturação de dados
\choice Independência de dados
\choice Controle de acesso e visualização de dados
\choice Controle de redundância e inconsistência
\choice Eficiência
\end{choices}

\question

A derivada da função \( f(x) = x^x \) é igual a:
\begin{choices}
\choice \( x x^{x-1} \)
\choice \( x^x \ln(x) \)
\choice \( x^x \)
\choice \( x^x (\ln(x) + 1) \)
\choice \( x^x (\ln(x) + x) \)
\end{choices}

\question

O conjunto básico de atividades e a ordem em que são realizadas no processo de construção de um software definem o que é habitualmente denominado de ciclo de vida do software. O ciclo de vida tradicional (também denominado waterfall) ainda é hoje em dia um dos mais difundidos e tem por característica principal: 
\begin{choices}
\choice a abordagem sistemática para realização das atividades do desenvolvimento de software de modo que elas seguem um fluxo sequencial; 
\choice a priorização da análise dos riscos do desenvolvimento
\choice a avaliação constante dos resultados intermediários feita pelo cliente.
\choice a codificação de uma versão executável do sistema desde as fases iniciais do desenvolvimento, de modo que o sistema final é incrementalmente construído, daí a alusão à ideia de "cascata" (waterfall); 
\choice o uso de formalização rigorosa em todas as etapas de desenvolvimento; 
\end{choices}



    %%%%%%%%%%%%%%%%%%%%%%%%%%%%%%%%%%%%%%%%%%%%%%%%%%%%%%%%%%%%%%%%%%%%%%%%%%%%%%%%
    %%% Finaliza o bloco de questões:
    \end{questions}
    \end{document}
    